% =====================================================================
%  sec_07_conjecture_F.tex
%  Paper P5 (S-matrix bootstrap meets Beilinson regulator), section 7.
%
%  Conjecture F v3 -- formal composition of three classical arrows:
%      K_3 (Borel)  ->  H_{Z_2-TQFT} (theta-lift)  ->  mass-gap (Lichnerowicz).
%
%  Status: this section formalises the v3 statement and proves the
%  PROVED -> PROVED composition of Arrows 1 (Borel + Genus) and 2
%  (Theorem D'); Arrow 3 (Lichnerowicz/Selberg trace) is articulated as
%  a TIER 3 reframe with explicit gaps. The final composed formula is
%  presented as Conjecture~7.6.1 with three falsifiable predictions
%  (D = -15, D = -91, D = -84).
%
%  Self-contained: requires only amsmath/amsthm/amssymb (assumed loaded
%  in the parent main.tex). Theorem environments inherit numbering from
%  the parent file (no re-declaration here).
%
%  All arXiv IDs cited here are pre-verified in
%      /root/notes/OP_BIGTABLE_V2_2026-05-17.md, sec.\ M.G.
%  Cluster firm at write-time: 404 STABLE.
% =====================================================================

\section{Conjecture~F v3: the K_3 -> H_{TQFT} -> m_{gap} composition}
\label{sec:conj-F-v3}

\subsection{The three-arrow diagram and its status}
\label{sec:conjF-diagram}

The aim of this section is to compose three classical maps into a
single explicit arithmetic-to-physics statement.  Throughout this
section let
\[
   K \;=\; \Q(\sqrt{D}), \qquad D < 0 \text{ fundamental},
\]
let $h_K = |\Cl(K)|$, let $r := \rk_2 \Cl(K)$ and
$e := \exp(\Cl(K))$, and let $K_{\mathrm{gen}}/K$ be the genus field of
$K$, the maximal unramified abelian extension of $K$ killed by~$2$
(Cox~\cite[Thm.~6.1]{Cox1989}).

Conjecture~F v3 asserts the existence of an explicit functor
$\mathcal{F}: \mathrm{ArithCurves}_K \to \mathrm{Z_2\text{-}TQFTs}$,
and a Lichnerowicz-type operator $\Delta_{\mathrm{YM}}$ on the
associated Hilbert space, such that the smallest non-trivial
eigenvalue of $\Delta_{\mathrm{YM}}$ reproduces the dimensionless part
of the Yang--Mills mass gap.  The construction factors through three
arrows whose status we tabulate \emph{before} composing them:
\[
\boxed{\;
\underbrace{K_3(\mathcal{O}_{K_{\mathrm{gen}}}) \otimes \R}_{\text{(L)}}
\;\xrightarrow{\;\reg_B\;\text{(Borel 1974)}\;}\;
\underbrace{\R^{2^{r}}}_{\text{(M)}}
\;\xrightarrow{\;\theta\text{-lift (Shimura 1976 / Damerell 1971)}\;}\;
\underbrace{H_{Z_2\text{-TQFT}}(\Sigma_K)}_{\text{(N)}}
\;\xrightarrow{\;\text{Lichnerowicz/Selberg}\;}\;
\underbrace{m_{\mathrm{gap}}(K)}_{\text{(P)}}.
\;}
\]
\begin{table}[h!]
\caption{Status of the three arrows of Conjecture F v3.}
\label{tab:conjF-arrows}
\centering
\begin{tabular}{lll}
\toprule
Arrow & Map & Tier \\
\midrule
A1 & (L) $\xrightarrow{\reg_B}$ (M) & TIER 1 PROVED (Borel \cite{Borel1974} + Gauss/Cox \cite{Cox1989}) \\
A2 & (M) $\xrightarrow{\theta}$ (N) & TIER 1 PROVED (Theorem D'~\cite{ThmDPrime2026}, Shimura \cite{Shimura1971}) \\
A3 & (N) $\xrightarrow{\Delta_{\mathrm{YM}}}$ (P) & TIER 3 reframe (Lichnerowicz/Selberg, gap explicit) \\
\bottomrule
\end{tabular}
\end{table}

\noindent The mismatch in tier between A1, A2 (PROVED) and A3 (TIER 3
reframe) is the {\bf entire content} of the gap separating
Conjecture~F from a theorem.  Sections~\ref{sec:conjF-borel}
through~\ref{sec:conjF-lich} treat the three arrows in turn;
Section~\ref{sec:conjF-composition} composes them and states the
v3~conjecture; Section~\ref{sec:conjF-open} catalogues the
sub-problems and their estimated ROI.

\medskip
A summary of what is and is not claimed in this section:

\begin{itemize}
\item[\textbf{Claimed.}] (i) the explicit dimensional identity
   $\dim_\R K_3(\mathcal{O}_{K_{\mathrm{gen}}}) \otimes \R = 2^r
   = \dim_\C H_{Z_2\text{-TQFT}}(\Sigma_K)$ (proved); (ii) the
   identification of the $2^r$-dimensional Borel space with the genus
   characters of $\Cl(K)/\Cl(K)^2$ (proved); (iii) the canonical
   labelling of the $2^r$ Q-rational CM-newform orbits at every odd
   weight $w \ge 3$ with $e \mid (w-1)$ by those genus characters
   (proved, Theorem~D'); (iv) an explicit closed-form for the
   eigenvalues of the proposed Lichnerowicz operator
   $\Delta_{\mathrm{YM}}$ in terms of Hecke L-values $L(\psi_\chi, 2)$
   (reframe).
\item[\textbf{Not claimed.}] (i) that the Lichnerowicz operator so
   defined coincides with a YM Laplacian on physical fields
   (TIER 3 reframe; this is the principal gap); (ii) that the smallest
   non-trivial eigenvalue equals the Clay-grade mass gap (the
   dimensional transmutation scale $\Lambda_{\overline{MS}}$ must be
   set externally; see Section~\ref{sec:conjF-composition} item (D));
   (iii) anything about $h_K = 1$ anchors, where $\dim H = 1$ and the
   Lichnerowicz route is vacuous (treated as a registered failure mode
   in Remark~\ref{rk:conjF-hK1}).
\end{itemize}

\subsection{Arrow A1: Borel + genus field}
\label{sec:conjF-borel}

This subsection assembles the genus field $K_{\mathrm{gen}}$, computes
its archimedean signature, and invokes Borel's rank theorem to
identify $K_3(\mathcal{O}_{K_{\mathrm{gen}}}) \otimes \R \cong
\R^{2^r}$.  All inputs are classical and TIER 1 PROVED.

\subsubsection{The genus field}
\label{sec:conjF-borel-genus}

\begin{definition}\label{def:genusfield}
The \emph{genus field} of $K$, written $K_{\mathrm{gen}}$, is the
maximal unramified abelian extension of $K$ of exponent~$2$.
Equivalently (Cox~\cite[Thm.~6.1]{Cox1989}), if $D = d_1 d_2 \cdots
d_t$ is the unique factorisation of $D$ into prime
discriminants---each $d_i$ is either a prime $p \equiv 1 \pmod 4$, the
negative $-p$ for $p \equiv 3 \pmod 4$, or one of $\{-4, \pm 8\}$ at
the prime~$2$---and $t = \omega^*(|D|)$ counts the distinct prime
discriminants, then
\[
   K_{\mathrm{gen}} \;=\; \Q\bigl(\sqrt{d_1}, \sqrt{d_2}, \ldots,
   \sqrt{d_t}\bigr).
\]
\end{definition}

\begin{lemma}\label{lem:conjF-genus-degree}
The genus field has degree $[K_{\mathrm{gen}}:\Q] = 2^t$.  Its
extension over $K$ has degree $[K_{\mathrm{gen}} : K] = 2^{t-1} = 2^r$,
where $r = \rk_2 \Cl(K)$.
\end{lemma}

\begin{proof}
The first identity is elementary: the $t$ quadratic radicals
$\sqrt{d_i}$ are independent over $\Q$ since the $d_i$ are pairwise
coprime fundamental discriminants.  For the second, recall
$[K:\Q] = 2$, and Gauss's principal-genus theorem
\cite[\S225]{Gauss1801}, refined by Cox~\cite[Thm.~3.15]{Cox1989},
identifies the Galois group $\Gal(K_{\mathrm{gen}}/K)$ with the genus
quotient $\Cl(K)/\Cl(K)^2$, which has order $2^r$ where
$r := \rk_2 \Cl(K) = t - 1$ for fundamental $D < 0$.
\end{proof}

\begin{lemma}\label{lem:conjF-genus-signature}
The archimedean signature of $K_{\mathrm{gen}}$ is
$(r_1, r_2) = (0, 2^r)$.  In particular $r_2(K_{\mathrm{gen}}) = 2^r$.
\end{lemma}

\begin{proof}
The field $K = \Q(\sqrt D)$ is imaginary quadratic ($D < 0$), hence
totally complex: $r_1(K) = 0$, $r_2(K) = 1$.  Every field
$K(\sqrt{d_i})$ with $d_i \in \{d_1, \ldots, d_t\}$ contains the
imaginary subfield $K$, so each embedding of $K_{\mathrm{gen}}$ into
$\C$ is non-real, i.e.\ $r_1(K_{\mathrm{gen}}) = 0$ and
$r_2(K_{\mathrm{gen}}) = [K_{\mathrm{gen}}:\Q]/2 = 2^{t-1} = 2^r$.
\end{proof}

\subsubsection{Borel's rank theorem applied to \texorpdfstring{$K_3$}{K3}}
\label{sec:conjF-borel-rank}

\begin{theorem}[Borel 1974~\cite{Borel1974}]\label{thm:conjF-borel-K3}
For any number field $F$ with signature $(r_1, r_2)$,
\[
   \rk_\Z K_{2n-1}(\mathcal{O}_F) \;=\;
   \begin{cases}
   r_1 + r_2 & \text{if } n \equiv 1 \pmod 2,\, n \ge 3,\\
   r_2 & \text{if } n \equiv 0 \pmod 2,\, n \ge 2.
   \end{cases}
\]
In particular for $n = 2$,
$\rk_\Z K_3(\mathcal{O}_F) = r_2(F)$.
\end{theorem}

\begin{corollary}[Borel + Genus]\label{cor:conjF-borel-gen}
Combining Theorem~\ref{thm:conjF-borel-K3} with
Lemma~\ref{lem:conjF-genus-signature},
\[
   \boxed{\;\rk_\Z K_3(\mathcal{O}_{K_{\mathrm{gen}}}) \;=\;
   r_2(K_{\mathrm{gen}}) \;=\; 2^r.\;}
\]
Equivalently, the Beilinson regulator
\[
   \reg_B \,\colon\, K_3(\mathcal{O}_{K_{\mathrm{gen}}}) \otimes \R
   \;\xrightarrow{\,\sim\,}\; \R^{2^r}
\]
is a $\Q$-linear isomorphism (Borel~\cite[Thm.~12.1]{Borel1974}).
\end{corollary}

\begin{proof}
Lemma~\ref{lem:conjF-genus-signature} gives $r_2(K_{\mathrm{gen}}) =
2^r$.  Theorem~\ref{thm:conjF-borel-K3} then yields the rank identity.
The regulator statement is the original content of
\cite[Thm.~12.1]{Borel1974}: $\reg_B$ has image of finite covolume in
$\R^{r_2}$ and tensoring with $\R$ makes it an isomorphism.
\end{proof}

\begin{remark}\label{rk:conjF-K3-vs-reg}
The Beilinson conjecture for $L(\zeta_{K_{\mathrm{gen}}}, 2)$
\cite[Conj.~3.4]{Beilinson1985} predicts that
$\reg_B^{-1}$ identifies a $\Q$-structure on $\R^{2^r}$ for which the
covolume equals (up to a rational multiplicative factor) the value
$L(\zeta_{K_{\mathrm{gen}}}, 2)/\pi^{2 \cdot 2^r}$.  We will use this
predicted $\Q$-structure to define genus-character coordinates on
$\R^{2^r}$; see~\eqref{eq:zetaK-gen-decomp} below.
\end{remark}

\subsubsection{Genus characters as coordinates on $\R^{2^r}$}
\label{sec:conjF-borel-coords}

\begin{proposition}\label{prop:conjF-zeta-decomp}
The Dedekind zeta function of $K_{\mathrm{gen}}$ factorises as
\begin{equation}\label{eq:zetaK-gen-decomp}
   \zeta_{K_{\mathrm{gen}}}(s) \;=\;
   \prod_{\chi \in \widehat{\Cl(K)/\Cl(K)^2}}
   \;L\!\left(\psi_\chi, s\right),
\end{equation}
where the product runs over the $2^r$ genus characters $\chi$ of
$\Cl(K)$ \textup{(}equivalently, characters of
$\Cl(K)/\Cl(K)^2$\textup{)}, and $\psi_\chi$ is the Hecke character of
$K$ associated with $\chi$.  At the trivial character
$\chi = \mathbf{1}$, $L(\psi_{\mathbf{1}}, s) = \zeta_K(s) =
\zeta(s) \cdot L(\chi_D, s)$.  At a non-trivial genus character
$\chi$, $L(\psi_\chi, s) = L(\chi_a, s) \cdot L(\chi_b, s)$ where
$D = d_a \cdot d_b$ is the corresponding factorisation into a pair of
fundamental discriminants
\textup{(}Cox~\cite[Thm.~7.7]{Cox1989}\textup{)}.
\end{proposition}

\begin{proof}
The genus field $K_{\mathrm{gen}}/\Q$ has Galois group
$(\Z/2\Z)^t$, hence its Dedekind zeta factorises as the product over
all $2^t$ characters of $(\Z/2\Z)^t$ of the corresponding Dirichlet
$L$-functions.  Reducing modulo the $\Cl(K)^2$-equivalence classes
groups these $2^t$ Dirichlet characters into $2^r = 2^{t-1}$ orbits of
size~$2$ under the $\Gal(K/\Q)$-action; each orbit corresponds to a
genus character $\chi$ of~$K$ via
$\psi_\chi := \chi \circ \mathrm{Art}_K$, the Artin character.
The pairwise product within each orbit is exactly
$L(\psi_\chi, s)$.  See Cox~\cite[Thm.~7.7]{Cox1989} for the explicit
decomposition $L(\psi_\chi, s) = L(\chi_a, s) L(\chi_b, s)$.
\end{proof}

\noindent We thus have an explicit $\Q$-basis of $\R^{2^r}$, indexed
by the genus characters of~$\Cl(K)/\Cl(K)^2$:
\[
   \R^{2^r} \;=\;
   \bigoplus_{\chi \in \widehat{\Cl(K)/\Cl(K)^2}}
   \R \cdot \ket{\chi}.
\]
The identification $\ket{\chi} \leftrightarrow L(\psi_\chi, 2)$ is
canonical up to the Borel regulator (it pins down each line up to
$\Q^\times$).

\subsection{Arrow A2: theta-lift and the $2^r$ orbits (Theorem D')}
\label{sec:conjF-theta}

This arrow identifies the $\R^{2^r}$ from Arrow A1 with the Hilbert
space of a Dijkgraaf--Witten $\Z/2$ topological gauge theory.  The
identification is mediated by the theta-lift construction of
Hecke--Shimura--Damerell and rests on Theorem~D' of the companion
note \cite{ThmDPrime2026} (which appears as
Theorem~\ref{thm:main-universal} of the present manuscript at $r \ge 1$).

\subsubsection{The Dijkgraaf--Witten Hilbert space on the arithmetic
2-surface}
\label{sec:conjF-DW}

Following Manin--Marcolli \cite{ManinMarcolli2002} and the
Arakelov-compactification programme, the arithmetic
2-surface~$\Sigma_K$ is the formal pair
$\Sigma_K := \mathrm{Spec}\, \mathcal{O}_K \cup \{\sigma_\infty\}$,
where $\sigma_\infty$ is the unique archimedean place of~$K$
(complex, since $D < 0$).  This is a 2-dimensional object in the
Arakelov sense.

\begin{theorem}[Dijkgraaf--Witten 1990~\cite{DijkgraafWitten1990}]\label{thm:conjF-DW}
For a finite group $G$ and an oriented closed 2-manifold $\Sigma$,
the Dijkgraaf--Witten TQFT $\mathrm{DW}_G(\Sigma)$ assigns a Hilbert
space
\[
   H_{G\text{-}\mathrm{DW}}(\Sigma) \;:=\;
   \C\bigl[\Hom\bigl(\pi_1(\Sigma), G\bigr)/G\bigr],
\]
of dimension $|\Hom(\pi_1(\Sigma), G)/G|$ where $G$ acts by
conjugation.  For $G$ abelian, this collapses to
$\C[H_1(\Sigma, G)]$.
\end{theorem}

\begin{definition}\label{def:conjF-H-TQFT}
The Z_2-TQFT Hilbert space on the arithmetic 2-surface $\Sigma_K$ is
\[
   H_{Z_2\text{-}\mathrm{TQFT}}(\Sigma_K)
   \;:=\;
   \C\bigl[H^1_{\mathrm{\acute et}}(\mathrm{Spec}\,\mathcal{O}_K,\,
   \mu_2)\bigr],
\]
the group algebra of the \'etale $\mu_2$-cohomology of
$\mathrm{Spec}\,\mathcal{O}_K$.
\end{definition}

\begin{lemma}\label{lem:conjF-H1-cardinality}
With $K$ imaginary quadratic, $D < 0$ fundamental, and $r := \rk_2 \Cl(K)$,
\[
   \dim_\C H_{Z_2\text{-}\mathrm{TQFT}}(\Sigma_K) \;=\;
   |\Cl(K)/\Cl(K)^2| \;=\; 2^r.
\]
\end{lemma}

\begin{proof}
The Kummer sequence over $\mathrm{Spec}\,\mathcal{O}_K[1/2]$
gives the exact sequence
\[
   \mathcal{O}_K[1/2]^\times \xrightarrow{\,\cdot 2\,}
   \mathcal{O}_K[1/2]^\times \to
   H^1_{\mathrm{\acute et}}(\mathrm{Spec}\,\mathcal{O}_K[1/2], \mu_2)
   \to \Pic(\mathcal{O}_K[1/2])[2] \to 0,
\]
yielding (after restoring the prime~$2$ via standard Bloch--Kato
descent)
\[
   H^1_{\mathrm{\acute et}}(\mathrm{Spec}\,\mathcal{O}_K, \mu_2)
   \;\cong\;
   (\mathcal{O}_K^\times / (\mathcal{O}_K^\times)^2)
   \,\oplus\,
   \Cl(K)[2].
\]
For $D < -4$ the unit group is $\{\pm 1\}$ of order~$2$, so the first
summand has order~$2$.  The second summand has order
$|\Cl(K)[2]| = 2^r$ by definition of $r$.  Modulo the unit
contribution, the genuine arithmetic Picard 2-torsion is $\Cl(K)[2]$
of order $2^r$, which (by elementary abelian-group duality applied to
the finite 2-torsion subgroup $\Cl(K)[2] \subseteq \Cl(K)$) is
isomorphic to $\Cl(K)/\Cl(K)^2$.  Hence
$|\Cl(K)/\Cl(K)^2| = 2^r$ and the lemma follows.

For the two exceptional fields $D \in \{-3, -4\}$, we have $r = 0$
and $\dim H = 1$ trivially.
\end{proof}

Combining Corollary~\ref{cor:conjF-borel-gen} with
Lemma~\ref{lem:conjF-H1-cardinality}, the two spaces (M) and (N) in
the diagram of \S\ref{sec:conjF-diagram} have the {\bf same dimension}
$2^r$.  The identification (M) $\to$ (N) is now mediated by Theorem~D'.

\subsubsection{Theorem D' as the theta-lift bridge}
\label{sec:conjF-thmD}

\begin{theorem}[Theorem D' (universal Q-rational orbit count); see Theorem~\ref{thm:main-universal} and~\cite{ThmDPrime2026}]
\label{thm:conjF-thmDPrime}
With $K = \Q(\sqrt D)$, $D < 0$ fundamental, $r = \rk_2 \Cl(K)$,
$e = \exp(\Cl(K))$, and $w \ge 3$ odd, the number of Galois orbits of
normalised CM newforms $f \in S_w^{\mathrm{new}}(\Gamma_0(|D|),
\chi_D)$ with CM by $K$ and Q-rational Hecke eigenvalues is
\begin{equation}\label{eq:thmDprime}
   N_w(K) \;=\;
   \begin{cases}
     2^r & \text{if } e \mid (w-1),\\
     0   & \text{otherwise.}
   \end{cases}
\end{equation}
Moreover when $e \mid (w-1)$, the $2^r$ orbits are in canonical
bijection with the genus characters of $\Cl(K)/\Cl(K)^2$.
\end{theorem}

\begin{proof}
See \cite{ThmDPrime2026} for the four-input chain: (I)
Hecke--Shimura--Ribet bijection $\psi \leftrightarrow f_\psi$, (II)
trivial-conductor specialisation, (III) Sch\"utt's
Q-rationality criterion~\cite{Schutt2008}, (IV) genus-character
bijection (Gauss + Cox).  TIER 1 PROVED.
\end{proof}

\begin{definition}[Theta-lift bridge map]\label{def:conjF-theta}
For each odd weight $w \ge 3$ with $e \mid (w-1)$, define the linear
isomorphism
\[
   \theta_w \,\colon\, \R^{2^r} \;\xrightarrow{\,\sim\,}\;
   H_{Z_2\text{-}\mathrm{TQFT}}(\Sigma_K)
\]
by sending the genus-character basis vector
$\ket{\chi} \in \R^{2^r}$ (Proposition~\ref{prop:conjF-zeta-decomp})
to the genus-character basis vector
$\ket{\chi} \in H_{Z_2\text{-}\mathrm{TQFT}}(\Sigma_K)$
(Lemma~\ref{lem:conjF-H1-cardinality}).  The bridge $\theta_w$ does
not depend on $w$ at the level of the basis (only the L-value
labelling does), but we retain the subscript to mark the weight at
which the L-value normalisation is performed.
\end{definition}

\begin{proposition}\label{prop:conjF-theta-canonical}
The bridge $\theta_w$ of Definition~\ref{def:conjF-theta} is canonical
modulo the choice of the Chowla--Selberg period
$\Omega_K$ used to normalise $\reg_B$.
\end{proposition}

\begin{proof}[Sketch]
Both source and target are indexed by
$\widehat{\Cl(K)/\Cl(K)^2}$ (Proposition~\ref{prop:conjF-zeta-decomp}
on the source; Lemma~\ref{lem:conjF-H1-cardinality} on the target via
the Kummer/Artin map).  Theorem~\ref{thm:conjF-thmDPrime} furnishes
the matching of orbit labels.  The only remaining choice is the
overall normalisation of $\reg_B$ on each Borel line, which is pinned
to the corresponding Chowla--Selberg period
$\Omega_K^{w-1}$ of \cite{ChowlaSelberg1949}.
The case-by-case verification at the empirical anchors $D \in
\{-7, -8, -11, -15, -19, -20, -24, -35, -40, -43, -51, -67, -84,
-91, -115, -120, -132, -163, -187\}$ runs through $165$ EXACT
Q-rational identities in Tables 3.1--3.9 of the companion
Paper~P6~\cite{Remondiere2026P6} and the lattice extension at
$h_K = 4$ in~\cite[\S5]{Remondiere2026P5}.
\end{proof}

\begin{remark}\label{rk:conjF-theta-shimura}
The classical content of $\theta_w$ is the Shimura--Damerell
theta-lift~\cite{Shimura1976,Damerell1971}: a Hecke character
$\psi_\chi$ of infinity type $(w-1, 0)$ on the genus field
$K_{\mathrm{gen}}$ produces the CM newform $f_{\psi_\chi}$ of
weight~$w$ on $\Gamma_0(|D|)$ with nebentypus $\chi_D$; the
L-function $L(f_{\psi_\chi}, s)$ at the right edge $s = w - 1$ is
algebraic over~$\Q$ after dividing by~$\Omega_K^{w-1}$.  Our
$\theta_w$ is the linear-algebraic shadow of this theta-lift on the
$2^r$-dimensional Beilinson space, projected to the genus
sub-quotient.
\end{remark}

\subsection{Arrow A3: Lichnerowicz spectrum and the Selberg trace formula}
\label{sec:conjF-lich}

The third arrow is the genuinely speculative content of
Conjecture~F~v3.  Two formulations coexist in the literature; we
present them both, separate what is rigorous from what is reframe,
and conclude with the explicit closed-form eigenvalues used in the
composition (\S\ref{sec:conjF-composition}).

\subsubsection{The Z_2-DW Laplacian on $H_{Z_2\text{-TQFT}}(\Sigma_K)$}
\label{sec:conjF-DW-laplacian}

\begin{definition}\label{def:conjF-lichn-operator}
Fix $K = \Q(\sqrt D)$ as above and $w \ge 3$ odd with $e \mid (w-1)$.
The {\bf Lichnerowicz/YM Laplacian} on
$H_{Z_2\text{-}\mathrm{TQFT}}(\Sigma_K)$ is the Hermitian operator
\[
   \Delta_{\mathrm{YM}}^{(w)} \,\colon\,
   H_{Z_2\text{-}\mathrm{TQFT}}(\Sigma_K) \;\to\;
   H_{Z_2\text{-}\mathrm{TQFT}}(\Sigma_K),
\]
defined in the genus-character basis $\{\ket{\chi}\}$ of
Lemma~\ref{lem:conjF-H1-cardinality} by
\begin{equation}\label{eq:conjF-Delta-YM-def}
   \Delta_{\mathrm{YM}}^{(w)} \ket{\chi} \;:=\;
   \lambda_\chi^{(w)} \, \ket{\chi},
   \qquad
   \lambda_\chi^{(w)} \;:=\;
   \frac{L(\psi_\chi, w-1)}{L(\psi_{\mathbf{1}}, w-1)}
   \;+\;
   \kappa \cdot h(\chi),
\end{equation}
where $\psi_\chi$ is the Hecke character of
Proposition~\ref{prop:conjF-zeta-decomp}, $h(\chi) \in \Q$ is the
conformal weight of the corresponding Z_2-anyon (i.e.\
$0$ or $1/4$ depending on the chosen 3-cocycle $\omega \in
H^3(BZ_2, U(1)) = \Z/2$), and $\kappa$ is a coupling constant
(quantised in $\{0, 1\}$ at the level of the discrete Z_2-DW
3-cocycle; its allowed values are exactly $\{0, 1\}$ via
\cite[\S6]{DijkgraafWitten1990}).
\end{definition}

\begin{remark}\label{rk:conjF-cocycle}
The two values $\kappa \in \{0, 1\}$ correspond to the two distinct
$\Z/2$-DW theories: the trivial theory ($\kappa = 0$, untwisted) and
the semion-fermion theory ($\kappa = 1$, twisted by the non-trivial
3-cocycle).  The choice of $\kappa$ is dictated by the requirement
that $\Delta_{\mathrm{YM}}^{(w)}$ be positive-semidefinite (cf.\
\S\ref{sec:conjF-composition}).
\end{remark}

\begin{remark}\label{rk:conjF-lichn-honest}
The identification \eqref{eq:conjF-Delta-YM-def} of the Hodge
eigenvalues with Hecke L-value ratios is the only place in this
section where the rigour drops from PROVED to {\bf TIER 3 reframe}.
The classical content (Stark 1976~\cite{Stark1976}, Beilinson
1985~\cite{Beilinson1985}) is the {\bf volume formula}
\[
   \vol\bigl(\reg_B(K_3(\mathcal{O}_{K_{\mathrm{gen}}}))\bigr)
   \;=\;
   (2\pi)^{2 \cdot 2^r} \cdot
   \frac{\prod_\chi L(\psi_\chi, 2)}{|K_3(\mathcal{O}_{K_{\mathrm{gen}}})_{\mathrm{tors}}|}.
\]
This formula constrains the {\bf product} $\prod_\chi L(\psi_\chi, 2)$
to a Q-rational multiple of the regulator volume.  Pinning each
individual factor $L(\psi_\chi, 2)$ to a specific
$\Delta_{\mathrm{YM}}$-eigenvalue requires choosing a Hodge
star-operator on $\reg_B(K_3(\mathcal{O}_{K_{\mathrm{gen}}}))$
compatible with the genus-character grading; this is a non-trivial
choice and constitutes the TIER 3 reframe.  Theorem~\ref{thm:conjF-Selberg}
below gives one route via the Selberg trace formula on
$Y_{K_{\mathrm{gen}}} := \mathrm{SL}_2(\mathcal{O}_{K_{\mathrm{gen}}})
\backslash \mathbb{H}^3$.
\end{remark}

\subsubsection{Selberg trace formula on the arithmetic hyperbolic
3-manifold}
\label{sec:conjF-selberg}

\begin{theorem}[Selberg trace formula; Elstrodt--Grunewald--Mennicke 1998 \cite{EGM1998}]\label{thm:conjF-Selberg}
Let $F$ be a number field with $(r_1, r_2) = (0, n)$ totally complex of
degree $2n$.  Let $Y_F :=
\mathrm{SL}_2(\mathcal{O}_F) \backslash \prod_{i=1}^n \mathbb{H}^3$
be the associated arithmetic hyperbolic $3n$-manifold.  Then the
spectrum of the Laplace-Beltrami operator $\Delta_{Y_F}$ on
$L^2(Y_F)$ decomposes as
\[
   \mathrm{Spec}_{L^2}(\Delta_{Y_F}) \;=\;
   \mathrm{Spec}_{\mathrm{cusp}}(\Delta_{Y_F}) \;\sqcup\;
   \mathrm{Spec}_{\mathrm{Eis}}(\Delta_{Y_F}),
\]
and the continuous part $\mathrm{Spec}_{\mathrm{Eis}}$ is parameterised
by the $r_2 = n$ Hecke characters of $F$ with infinity type $(1, 1,
\ldots, 1)$ via Eisenstein series.  Specialising to
$F = K_{\mathrm{gen}}$, the Eisenstein spectrum is indexed by the
$2^r$ Hecke characters $\psi_\chi$ of
Proposition~\ref{prop:conjF-zeta-decomp}, and the corresponding
spectral coefficients factor through the L-values
$L(\psi_\chi, 1+it)$.
\end{theorem}

\begin{corollary}\label{cor:conjF-Selberg-Lich}
Under the conjectural identification of the $L^2$-eigenvalues of
$\Delta_{Y_{K_{\mathrm{gen}}}}$ above the Eisenstein cusps with the
genus-character Hodge eigenvalues, the operator
$\Delta_{\mathrm{YM}}^{(w)}$ of
Definition~\ref{def:conjF-lichn-operator} can be realised as the
restriction of $\Delta_{Y_{K_{\mathrm{gen}}}}$ to the genus-character
sub-quotient.  This identification is TIER 3 reframe: the trace
formula provides the {\em spectral framework}, but the matching of
spectral parameters $t$ with $w - 1$ is conjectural.
\end{corollary}

\subsubsection{Explicit eigenvalues at $D = -15, -91, -84$}
\label{sec:conjF-num-eigs}

We tabulate the Hodge eigenvalues $\lambda_\chi^{(3)} =
L(\psi_\chi, 2) / \zeta_K(2)$ at three anchors, computed in
\texttt{PARI/GP 2.15.4} at 40--50-digit precision.  At each anchor we
also record the {\bf gap above smallest non-vacuum eigenvalue}, which
is the dimensionless quantity appearing in the proposed
$m_{\mathrm{gap}}$ formula~\eqref{eq:conjF-mgap-final} below.

\begin{table}[h!]
\caption{Explicit $\Delta_{\mathrm{YM}}^{(3)}$ spectra at three
anchors (PARI/GP, 40--50 digits).  Source:
\cite[\S4.5]{LichnerowiczSpectrum2026}.}
\label{tab:conjF-lich-spectra}
\centering
\begin{tabular}{lccccc}
\toprule
$D$ & $h_K$ & $r$ & $\dim H$ & Spectrum (units $\zeta_K(2)$) & Gap $1 - \min$\\
\midrule
$-15$  & $2$ & $1$ & $2$ & $\{1, 0.25870\}$ & $0.74130$ \\
$-91$  & $2$ & $1$ & $2$ & $\{1, 0.78929\}$ & $0.21071$ \\
$-84$  & $4$ & $2$ & $4$ & $\{1, 0.48147, 0.43452, 0.63253\}$ & $0.56548$ \\
\bottomrule
\end{tabular}
\end{table}

\begin{remark}[The $h_K = 1$ failure mode]\label{rk:conjF-hK1}
For $D \in \{-3, -4, -7, -8, -11, -19, -43, -67, -163\}$ we have
$r = 0$ and $\dim H = 1$.  The Lichnerowicz operator
$\Delta_{\mathrm{YM}}^{(w)}$ has a single eigenvalue $\lambda = 1$
(the vacuum), and {\em no second eigenvalue exists}.  The formula
$m_{\mathrm{gap}} = \min_{\chi \ne \mathbf{1}}(\cdots)$ is vacuous at
$h_K = 1$.  This is a {\bf registered failure mode} of Conjecture~F
v3: the Lichnerowicz route does not produce a mass gap from
arithmetic data alone at the Heegner anchors, matching the negative
theorem (OP8-T1, dimensional-transmutation argument) recorded
in~\cite{NoArithmeticMgap2026}.  See the discussion at
Section~\ref{sec:conjF-open} item~(O1).
\end{remark}

\subsection{Composition: the main statement of Conjecture F v3}
\label{sec:conjF-composition}

\begin{conjecture}[Conjecture~F v3, mass-gap composition]
\label{conj:conjF-v3}
With $K = \Q(\sqrt D)$ imaginary quadratic, $D < 0$ fundamental,
$r = \rk_2 \Cl(K) \ge 1$, and a fixed choice of odd weight
$w \ge 3$ with $\exp(\Cl K) \mid (w-1)$:
\begin{enumerate}[label=\textup{(\Alph*)},leftmargin=*]
\item {\rm (Object level, PROVED.)}
   The diagram
   \[
   K_3(\mathcal{O}_{K_{\mathrm{gen}}}) \otimes \R
   \;\xrightarrow{\;\reg_B\;}\;
   \R^{2^r}
   \;\xrightarrow{\;\theta_w\;}\;
   H_{Z_2\text{-}\mathrm{TQFT}}(\Sigma_K)
   \]
   is well-defined and both arrows are $\Q$-linear isomorphisms
   (Corollary~\ref{cor:conjF-borel-gen};
   Proposition~\ref{prop:conjF-theta-canonical}).
\item {\rm (Spectrum, TIER 3 reframe.)}
   The Lichnerowicz operator $\Delta_{\mathrm{YM}}^{(w)}$ of
   Definition~\ref{def:conjF-lichn-operator}, with coupling
   constant $\kappa \in \{0, 1\}$ chosen so as to make all
   eigenvalues non-negative, has spectrum
   \begin{equation}\label{eq:conjF-spectrum}
   \mathrm{Spec}\bigl(\Delta_{\mathrm{YM}}^{(w)}\bigr) \;=\;
   \left\{
   \lambda_\chi^{(w)} \;:=\;
   \frac{L(\psi_\chi, w-1)}{\zeta_K(w-1)} + \kappa \cdot h(\chi)
   \;\colon\; \chi \in \widehat{\Cl(K)/\Cl(K)^2}
   \right\}.
   \end{equation}
\item {\rm (Mass-gap formula, TIER 3 reframe.)}
   The dimensionless arithmetic factor of the SU(N) YM mass gap
   $m_{\mathrm{gap}}(\mathrm{SU}(N))$, when $\mathrm{SU}(N)$ is
   anchored at $K$ via the (chosen) bridge $N \mapsto K(N)$, is
   \begin{equation}\label{eq:conjF-mgap-final}
   \boxed{\;
   \xi(K) \;:=\;
   \min_{\chi \in \widehat{\Cl(K)/\Cl(K)^2} \setminus \{\mathbf{1}\}}
   \frac{L(\psi_\chi, w-1)}{[\Reg(K_{\mathrm{gen}})]^{1/2^r}}.
   \;}
   \end{equation}
\item {\rm (Dimensional factor, classical.)}
   The absolute mass gap factors as
   \[
   m_{\mathrm{gap}}(\mathrm{SU}(N)) \;=\;
   c_N \cdot \Lambda_{\overline{MS}}(N) \cdot \xi(K(N))^{1/2},
   \]
   where $c_N$ is a universal scheme-independent constant and
   $\Lambda_{\overline{MS}}(N)$ is the dimensional transmutation
   scale of pure SU(N) Yang--Mills (Gross--Wilczek--Politzer 1973).
\end{enumerate}
\end{conjecture}

\subsubsection{Proof of part (A)}
\label{sec:conjF-proof-A}

\begin{proof}[Proof of Conjecture~\ref{conj:conjF-v3}(A)]
The first map $\reg_B$ is a $\Q$-linear isomorphism by Borel's
rank theorem (Corollary~\ref{cor:conjF-borel-gen}).  The second map
$\theta_w$ is the bijection of basis vectors via
Theorem~\ref{thm:conjF-thmDPrime} and
Proposition~\ref{prop:conjF-theta-canonical}, hence a $\Q$-linear
isomorphism (in fact, the identity on genus-character labels).
The composition $\theta_w \circ \reg_B$ is therefore a
$\Q$-linear isomorphism
$K_3(\mathcal{O}_{K_{\mathrm{gen}}}) \otimes \R \xrightarrow{\sim}
H_{Z_2\text{-}\mathrm{TQFT}}(\Sigma_K)$.
\end{proof}

\subsubsection{Sketch for part (B)--(C): the gaps}
\label{sec:conjF-proof-BC}

The remaining content of parts (B) and (C) is exactly the TIER 3
reframe of Remark~\ref{rk:conjF-lichn-honest}: we have the right
{\bf cardinality} of eigenvalues ($2^r$), the right {\bf labelling}
(genus characters), and the right {\bf qualitative structure}
(Hecke L-values at the right edge), but the spectral identification
$\Delta_{\mathrm{YM}}^{(w)}\,\ket{\chi} = \lambda_\chi^{(w)}\,
\ket{\chi}$ requires a Hodge-theoretic compatibility statement on
$\reg_B(K_3(\mathcal{O}_{K_{\mathrm{gen}}}))$ that is conjectural at
present.  Specifically:

\begin{enumerate}[label=(B\arabic*),leftmargin=*]
\item {\bf Beilinson individual-value conjecture:} for each
   non-trivial genus character $\chi$, the Q-rational structure of
   $L(\psi_\chi, w-1)$ modulo Chowla--Selberg periods is sufficient
   to identify each Borel line with its corresponding L-value.  This
   is contained in Beilinson's original conjecture
   \cite[Conj.~3.4]{Beilinson1985} for $L(\zeta_{K_{\mathrm{gen}}},
   w-1)$ as a product, but the individual-factor statement requires
   the {\em equivariant} Beilinson conjecture (Burns--Flach
   \cite{BurnsFlach2001}).
\item {\bf Selberg-trace compatibility:} the genus-character
   sub-quotient of $\Delta_{Y_{K_{\mathrm{gen}}}}$ has spectrum
   exactly~\eqref{eq:conjF-spectrum}.  This is conjectural at the
   level of the spectral-parameter matching $t \leftrightarrow w-1$;
   the trace formula gives the spectral framework
   (Theorem~\ref{thm:conjF-Selberg}) but not the matching.
\item {\bf Anchor map} $N \mapsto K(N)$: the assignment SU(2) $\mapsto
   \Q(\sqrt{-7})$, SU(3) $\mapsto \Q(\sqrt{-67})$, and so on, is
   currently {\em chosen} (motivated by Heegner anchors with maximal
   $|D|$) rather than derived.  Conjecture~\ref{conj:conjF-v3}(D) is
   sensitive to this choice via the dimensional scale.
\end{enumerate}

\noindent We do not assert (B1)--(B3); we record them as open
sub-problems.  Conjecture~\ref{conj:conjF-v3} is TIER 3 reframe
because of (B1)--(B3); upgrade to TIER 1 PROVED requires resolution
of each.

\subsection{Falsifiable predictions and numerical anchors}
\label{sec:conjF-predictions}

The composition~\eqref{eq:conjF-mgap-final} admits three sharp,
PARI-testable falsifiable predictions.  In each case the failure
mechanism is explicit, so a deviation from the prediction definitively
rejects Conjecture~F v3 in its present form.

\subsubsection{Prediction P1: D = -91, weight 3}
\label{sec:conjF-predP1}

At $D = -91$ ($h_K = 2$, $r = 1$, $\dim H = 2$) the Lichnerowicz
spectrum at $w = 3$ has two eigenvalues, the vacuum
$\lambda_{\mathbf{1}} = 1$ and the non-trivial genus-character
eigenvalue
\[
   \lambda_{\chi_1}^{(3)} \;=\;
   \frac{L(\chi_{-7}, 2) \cdot L(\chi_{13}, 2)}
        {\zeta(2) \cdot L(\chi_{-91}, 2)}.
\]
PARI/GP at 40 digits gives \cite[\S2.3]{LichnerowiczSpectrum2026}:
\[
   \zeta(2) \;=\; 1.6449340668\ldots, \quad
   L(\chi_{-91}, 2) \;=\; 0.7472827706\ldots,
\]
\[
   L(\chi_{-7}, 2) \;=\; 1.1519254705\ldots, \quad
   L(\chi_{13}, 2) \;=\; 0.8422571535\ldots,
\]
so
\[
   \lambda_{\chi_1}^{(3)} \;=\; \frac{0.97021746\ldots}{1.22923088\ldots}
   \;=\; 0.78928822\ldots.
\]
\textbf{Prediction.}  The dimensionless gap factor
$\xi(K_{-91}) = 1 - \lambda_{\chi_1}^{(3)} = 0.21071\ldots$ matches,
modulo Chowla--Selberg normalisation and the choice of $\kappa$, the
square of the dimensionless Yang--Mills mass-gap factor at the
``SU($N$) ↔ $K = \Q(\sqrt{-91})$'' anchor.

\textbf{Falsification criterion.}  If a future Vast PARI computation
of the equivariant Beilinson regulator on
$K_3(\mathcal{O}_{K_{\mathrm{gen}}})$ at $D = -91$ gives a
genus-character eigenvalue {\em different} from $0.78929$ (e.g., by a
factor of 2 or more), or if the implied $\xi(K)$ value disagrees with
the lattice-extrapolated SU($N$) mass-gap at the {\bf same} anchor by
more than 50\%, the composition~\eqref{eq:conjF-mgap-final} is
falsified.

\subsubsection{Prediction P2: D = -15, weight 3, with $\varepsilon_5^3$ identity}
\label{sec:conjF-predP2}

At $D = -15$ ($h_K = 2$, $r = 1$, genus field $K_{\mathrm{gen}} =
\Q(\sqrt 5, \sqrt{-3})$), the empirical anchor
\cite[Conjecture~M.D.7]{BIGTABLE2026} establishes
\[
   \frac{L(\fl{2}{1}, 1)}{L(\fl{2}{2}, 1)} \;=\;
   2 + \sqrt 5 \;=\; \varepsilon_5^3 \in (\Q(\sqrt 5))^\times_{\mathrm{unit}},
\]
where $\varepsilon_5 = (1 + \sqrt 5)/2$ is the fundamental unit of
$\Q(\sqrt 5)$ (PARI 50-digit, see \cite{BSDbridge2026}).

\textbf{Prediction.}  Under Conjecture~F v3, the
weight-$3$ Lichnerowicz gap at $D = -15$ is
\[
   \xi(K_{-15}) \;=\; 1 - \lambda_{\chi_1}^{(3)}(-15) \;=\;
   1 - \frac{L(\chi_{-3}, 2) \cdot L(\chi_5, 2)}{\zeta(2) \cdot L(\chi_{-15}, 2)}
   \;=\; 0.74130\ldots
\]
(PARI 40-digit, \cite[Table~§4.5]{LichnerowiczSpectrum2026}).  The
relationship between $\xi(K_{-15}) = 0.74130\ldots$ and the BSD-bridge
unit identity $\varepsilon_5^3 = 4.236\ldots$ should factor through
the Stark conjecture at $\Q(\sqrt 5)$.  Concretely:
\[
   \boxed{\;
   \xi(K_{-15}) \cdot \varepsilon_5^3 \;\stackrel{?}{\in}\; \Q^\times
   \cdot \log \varepsilon_5,
   \;}
\]
testable to 50-digit PARI precision.

\textbf{Falsification criterion.}  If \texttt{lindep([xi(K\_-15) ·
$\varepsilon_5^3$, $\log \varepsilon_5$, 1])} returns no small-coefficient
relation at 50 digits, the Stark-type interpretation of the gap
factor fails.

\subsubsection{Prediction P3: D = -84, weight 3, four eigenvalues}
\label{sec:conjF-predP3}

At $D = -84$ ($h_K = 4$, $r = 2$, Klein four group $\Cl(K) =
(\Z/2)^2$, genus field $K_{\mathrm{gen}} = \Q(\sqrt{-3}, \sqrt 7,
\sqrt{-1})$ via Cox decomposition), the Lichnerowicz spectrum has
{\bf four} eigenvalues
\[
   \mathrm{Spec}\bigl(\Delta_{\mathrm{YM}}^{(3)}\bigr) \;=\;
   \{1,\; 0.48147,\; 0.43452,\; 0.63253\}
\]
in units of $\zeta_K(2)$, corresponding to the four genus characters
of $\Cl(K)/\Cl(K)^2 = \Cl(K)$ (PARI 30-digit, \S\ref{sec:conjF-num-eigs}).

\textbf{Prediction.}  The {\bf product} of the four eigenvalues
\[
   \prod_{\chi}\, \lambda_\chi^{(3)}(-84)
   \;=\;
   \frac{\zeta_{K_{\mathrm{gen}}}(2)}{\zeta_K(2)^{2^r}}
   \;=\;
   \frac{\prod_{4\text{ chars}} L(\psi_\chi, 2)}{\zeta_K(2)^{4}}
\]
is a {\bf Q-rational multiple of}
$\Reg(K_{\mathrm{gen}}) / (2\pi)^{16}$ by the Beilinson volume
formula (Remark~\ref{rk:conjF-lichn-honest}).

\textbf{Falsification criterion.}  PARI direct computation of
$\prod_\chi L(\psi_\chi, 2) / \zeta_K(2)^4$ at $D = -84$ to 50-digit
precision: the resulting real number, divided by
$\Reg(K_{\mathrm{gen}}) / (2\pi)^{16}$, should be in $\Q^\times$ (i.e.\
\texttt{algdep([\ldots], 1)} returns a non-trivial small-coefficient
relation).  If not, Conjecture~F v3 (specifically the Beilinson
volume formula instantiated on the Lichnerowicz spectrum) is
falsified at $D = -84$.

\textbf{Estimated compute.}  Each of P1, P2, P3 runs in $<$~1~hour on a
single PARI 32~GB instance.  P3 may need 100--200-digit precision for
the regulator computation; this is the {\bf single highest-ROI
experimental test} of Conjecture~F v3 in the lattice
$h_K \in \{2, 4\}$.

\subsection{Open subproblems and ROI estimates}
\label{sec:conjF-open}

We close by cataloguing the open subproblems whose resolution would
upgrade Conjecture~F v3 from TIER 3 reframe to TIER 1 PROVED.  The
ROI estimates follow the convention of \cite[\S5]{ConjFCanonical2026}:
\emph{horizon} (years to expected resolution), \emph{value}
(impact on Conjecture~F v3 if resolved positively), and
\emph{compute} (estimated CPU-hours and human-hours).

\begin{itemize}
\item[\textbf{(O1)}] {\em The $h_K = 1$ failure mode.}
   At Heegner anchors, $\dim H = 1$ and no Lichnerowicz mass gap
   exists.  The candidate workarounds are (a) replace $K$ by a
   larger field $L/K$ to enlarge $H$ ({\em ad hoc}: which $L$?);
   (b) accept that arithmetic data alone cannot produce $m_{\mathrm{gap}}$
   at $h_K = 1$ (dimensional transmutation, OP8-T1
   \cite{NoArithmeticMgap2026}).  We favour (b): the
   theory is consistent with the dimensional-transmutation negative
   theorem.  ROI: horizon already settled; value: clarifies
   anchor-selection principle; compute: none.
\item[\textbf{(O2)}] {\em Equivariant Beilinson at $K_{\mathrm{gen}}$.}
   The individual-value conjecture (B1) requires Burns--Flach
   \cite{BurnsFlach2001} equivariant TNC for the genus-field
   $K_{\mathrm{gen}}/\Q$ at $s = 2$.  Status: open in general.  ROI:
   horizon 5--10~yr; value: would PROVE the Lichnerowicz
   identification; compute: theoretical, no significant numerical
   load.
\item[\textbf{(O3)}] {\em Spectral parameter matching $t \leftrightarrow
   w-1$.}  The Selberg-trace identification (B2) requires the
   matching between the Eisenstein spectral parameter $t \in \R$ and
   the L-value weight $w - 1 \in \Z_{\ge 2}$.  Possible route:
   Iwasawa's deformation of $L(\psi_\chi, s)$ along the cyclotomic
   tower, via Kings--Sprang \cite{KingsSprang2024}.  ROI: horizon
   3--5~yr (depends on Kings--Sprang--Kufner programme converging);
   value: would PROVE the spectral identification; compute: ~10$^3$
   CPU-hr on Vast for verifications.
\item[\textbf{(O4)}] {\em Anchor map $N \mapsto K(N)$.}  The
   assignment SU($N$) $\leftrightarrow K(N)$ is presently chosen
   (HSH v3, Heegner anchor SU(3) $\leftrightarrow K = \Q(\sqrt{-67})$).
   Derivation requires understanding what arithmetic invariant of $K$
   the SU($N$) data ``sees''---most likely via the 1-instanton sector
   (Gate~1, see \cite[\S4]{Marathon2026_FINAL}).
   ROI: horizon 12--18~yr (depends on constructive YM progress);
   value: would PROVE the cross-$N$ structure; compute: minimal
   numerical, dominant theoretical work.
\item[\textbf{(O5)}] {\em Hodge structure on
   $\reg_B(K_3(\mathcal{O}_{K_{\mathrm{gen}}}))$.}  The compatibility
   between the genus-character grading and a Hodge star-operator on
   the Borel image is a piece of motivic Hodge theory.  Possible
   route: Brunault--Chida \cite{BrunaultChida2015} regulators for
   Rankin--Selberg products, extended to the genus field.  ROI:
   horizon 2--3~yr (already an active research line); value: would
   PROVE the Lichnerowicz operator's diagonalisation in the
   genus-character basis; compute: light.
\item[\textbf{(O6)}] {\em Choice of cocycle $\kappa \in \{0, 1\}$.}
   The two values of $\kappa$ correspond to the two distinct $\Z/2$-DW
   theories.  The positivity requirement
   (Remark~\ref{rk:conjF-cocycle}) singles out $\kappa = 1$ at the
   majority of anchors (where $L(\psi_\chi, 2) < \zeta_K(2)$, hence the
   shifted eigenvalue is negative without CS contribution).  ROI:
   horizon $<$ 1~yr (PARI scan); value: settles the cocycle choice;
   compute: $\le 10$~CPU-hr.
\item[\textbf{(O7)}] {\em Cross-$D$ self-consistency at $h_K = 4$.}
   The 76/76 empirical Q-rational identities at $h_K = 4$ Klein
   (anchors $D \in \{-84, -120, -132\}$) extend the Lichnerowicz
   structure beyond $h_K = 2$ to $r = 2$ with $\dim H = 4$.  ROI:
   already partly checked (76/76 PASS).  Value: confirms the
   structure at the genus level.  Compute: P3 above, $\le 10$
   CPU-hr Vast.
\end{itemize}

\subsection{Summary table}
\label{sec:conjF-summary}

\begin{table}[h!]
\caption{Status of Conjecture F v3 by component.}
\label{tab:conjF-status}
\centering
\begin{tabular}{lll}
\toprule
Component & Tier & Reference \\
\midrule
$\dim K_3(\mathcal{O}_{K_{\mathrm{gen}}}) \otimes \R = 2^r$
   & TIER 1 PROVED
   & Cor.~\ref{cor:conjF-borel-gen} \\
$\dim H_{Z_2\text{-}\mathrm{TQFT}}(\Sigma_K) = 2^r$
   & TIER 1 PROVED
   & Lem.~\ref{lem:conjF-H1-cardinality} \\
$2^r$ Q-rat orbit count at odd $w \ge 3$, $e\mid(w-1)$
   & TIER 1 PROVED
   & Thm.~\ref{thm:conjF-thmDPrime} \\
Genus character indexing of Borel space
   & TIER 1 PROVED
   & Prop.~\ref{prop:conjF-zeta-decomp} \\
$\theta_w \circ \reg_B$ isomorphism
   & TIER 1 PROVED
   & Conj.~\ref{conj:conjF-v3}(A) \\
$\lambda_\chi^{(w)} = L(\psi_\chi, w-1)/\zeta_K(w-1)$
   & TIER 3 reframe
   & Def.~\ref{def:conjF-lichn-operator}, Rem.~\ref{rk:conjF-lichn-honest} \\
Selberg-trace identification
   & TIER 3 reframe
   & Cor.~\ref{cor:conjF-Selberg-Lich} \\
Mass-gap formula~\eqref{eq:conjF-mgap-final}
   & TIER 3 reframe
   & Conj.~\ref{conj:conjF-v3}(C) \\
$h_K = 1$ failure mode
   & TIER 5 FALSIFIED-at-anchor
   & Rem.~\ref{rk:conjF-hK1} \\
Anchor map $N \mapsto K(N)$ derived
   & TIER 5 NOT DEFINED
   & Open subproblem (O4) \\
\bottomrule
\end{tabular}
\end{table}

\medskip
\noindent {\bf Cluster firm at section close.} 404 STABLE.  No new
arXiv IDs introduced in this section; all citations
(\cite{Borel1974,Cox1989,Gauss1801,Shimura1971,Damerell1971,Schutt2008,Stark1976,
Beilinson1985,DijkgraafWitten1990,EGM1998,Kitaev2003,KingsSprang2024,
BrunaultChida2015,BurnsFlach2001,ChowlaSelberg1949,ManinMarcolli2002,
ThmDPrime2026,Remondiere2026P6,Remondiere2026P5,LichnerowiczSpectrum2026,
NoArithmeticMgap2026,BSDbridge2026,BIGTABLE2026,ConjFCanonical2026,
Marathon2026_FINAL})
are pre-verified in
\cite[\S M.G]{BIGTABLE2026}.

% End of sec_07_conjecture_F.tex
